% TOP_PROBLEM: {"id": 30007253, "title": "Quaternionic division intersections of Hodge-generic curves in A2", "category_id": 1, "category": {"id": 1, "name": "number_theory", "display_name": "Number Theory", "description": "Properties of integers, prime numbers, Diophantine equations.", "slug": "number-theory", "order_index": 1, "created_at": "2026-07-31T15:26:25.670Z"}, "status": "open"}
% FIRST_POSTED: 2026-10-06
% ENTRY_KIND: statement_only
% !TeX program = lualatex
\documentclass[11pt]{article}
\usepackage[margin=25mm]{geometry}
\usepackage{fontspec}
\setmainfont{Latin Modern Roman}
\usepackage{amsmath,amssymb,mathtools}
\usepackage{unicode-math}
\setmathfont{Latin Modern Math}
\usepackage{xurl,hyperref}
\hypersetup{colorlinks=true,urlcolor=blue}
\setcounter{secnumdepth}{0}
\setlength{\parindent}{0pt}
\setlength{\parskip}{5pt}
\setlength{\emergencystretch}{3em}
\newcommand{\R}{\mathbb R}
\newcommand{\N}{\mathbb N}
\newcommand{\Z}{\mathbb Z}
\newcommand{\Q}{\mathbb Q}
\newcommand{\C}{\mathbb C}
\newcommand{\D}{\mathbb D}
\newcommand{\F}{\mathbb F}
\newcommand{\sref}[2]{\hyperlink{source:#1:#2}{[S#2]}}
\newcommand{\cataloguescope}[1]{#1}
\begin{document}
\section*{Quaternionic division intersections of Hodge-generic curves in A2}\label{30007253}
\subsection{Definitions and mathematical statement}
\paragraph{Statement recorded in the supplied source.}
For a reduced irreducible closed Hodge-generic Qbar-curve C in A2, are there only finitely many points whose abelian surface has rational endomorphism algebra an indefinite quaternion division algebra over Q?
\subsection{Short English statement}
For a reduced irreducible closed Hodge-generic Qbar-curve C in A2, are there only finitely many points whose abelian surface has rational endomorphism algebra an indefinite quaternion division algebra over Q?
\subsection{Sources}
\begin{itemize}
\item[S1] Original question source or mathematical context.
\url{https://en.wikipedia.org/w/index.php?title=List_of_unsolved_problems_in_mathematics&oldid=1366636767#Number_theory}.
\item[S2] Original question source or mathematical context.
\url{https://ems.press/content/serial-article-files/52570}.
\item[S3] Original question source or mathematical context.
\url{https://arxiv.org/html/2504.00865v1}.
\item[S4] Original question source or mathematical context.
\url{https://arxiv.org/html/2209.07967}.
\item[S5] OpenAI, manuscript citation page. Source of the formulation motivating this local problem.
\url{https://github.com/openai/math/blob/adc7f1241b42e322a6451854ab7e4b4c146bf78a/preprints/Quaternionic-Division-Points-on-Curves-in-the-Siegel-Threefold-September-24-2026/README.md}.
\end{itemize}
\end{document}
